\section*{Appendix C: Implementation Details}

In practice, we implement our estimator in PyTorch by leveraging automatic differentiation to compute both the standard autoencoder (AE) gradient and our Fisher-consistent estimating equation gradient, and combine them using a convex interpolation.

Recall that our estimating equation is defined as:
\[
\frac{1}{n} \sum_{i=1}^n \psi_\bt(\Y_i, e(\Y_i, \bt)) 
=
\mathbb{E}_{\Y \sim f_\bt} \left[ \psi_\bt(\Y, e(\Y, \bt)) \right],
\]
where \( \psi_\bt(y, z) = T(y) \cdot \partial \eta(z, \bt)/\partial \bt \) is our score-like function, and \( \eta(z, \bt) \) is the decoder's linear predictor. The encoder \( z = e_\phi(y) \) is implemented as a differentiable neural network.

To ensure stable training while maintaining statistical consistency, we define a hybrid estimating equation as a convex combination of the autoencoder's empirical moment and our model-based expectation:
\[
\text{LHS} = \alpha \cdot \text{RHS}_{\text{AE}} + (1 - \alpha) \cdot \text{RHS}_{\text{SME}},
\]
where:
\begin{itemize}
  \item \( \text{RHS}_{\text{AE}} = \frac{1}{n} \sum_{i=1}^n \psi_\bt(\Y_i, e_\phi(\Y_i)) \) is the usual autoencoder gradient term,
  \item \( \text{RHS}_{\text{SME}} = \mathbb{E}_{\Y^\star \sim f_\bt} [\psi_\bt(\Y^\star, e_\phi(\Y^\star))] \) is estimated using Monte Carlo samples from the model.
\end{itemize}

This leads to the decoder update:
\[
\nabla_\bt \mathcal{L} = \alpha \cdot \nabla_\bt \mathcal{L}_{\text{AE}} + (1 - \alpha) \cdot \nabla_\bt \mathcal{L}_{\text{SME}}.
\]

\paragraph{Implementation Details.}
To compute this in PyTorch, we proceed as follows:
\begin{enumerate}
    \item Perform a forward pass through the encoder and decoder to obtain reconstructed outputs and compute the AE loss.
    \item Backpropagate the AE loss using \verb|loss.backward(retain_graph=True)|. This populates the gradients for both the encoder (\( \phi \)) and decoder (\( \theta \)).
    \item Clone and store the decoder gradients (i.e., \verb|p.grad.clone()| for each \verb|p in decoder.parameters()|) as the AE component.
    \item Using \verb|torch.no_grad()|, sample synthetic data \( \Y^\star \sim f_\theta \) from the decoder, and encode these samples with \verb|e_\phi(\Y^\star)| without tracking gradients.
    \item Compute \( \psi_\bt(\Y_i, e_\phi(\Y_i)) \) and \( \psi_\bt(\Y^\star, e_\phi(\Y^\star)) \), and define the SME loss as their (batchwise) difference.
    \item Backpropagate the SME loss to accumulate the SME gradient into the decoder’s parameter \verb|.grad|.
    \item Linearly interpolate the stored AE and SME decoder gradients using \( \alpha \), and replace the decoder gradients with the weighted sum.
    \item Update decoder and encoder parameters using their respective optimizers. Only the AE gradient is used to update the encoder.
\end{enumerate}

This approach enables the encoder to be trained using stable AE-style gradients, while the decoder is progressively corrected toward the Fisher-consistent solution. The interpolation parameter \( \alpha \in [0, 1] \) can be annealed during training to start with AE-like behavior and converge to the statistically principled estimator.
